% =============================================================================
%  Orbital CDN: A Physically-Grounded, Predictive Control Plane for
%  Content Delivery over Low-Earth-Orbit Constellations
%
%  Compiles with a base TeX Live installation:
%     pdflatex orbital_cdn.tex   (run twice for cross-references)
%  No packages beyond the LaTeX base + geometry/booktabs/hyperref set.
% =============================================================================

\documentclass[10pt,twocolumn]{article}

\usepackage[a4paper,margin=1.9cm,columnsep=0.7cm]{geometry}
\usepackage{amsmath,amssymb}
\usepackage{booktabs}
\usepackage{tabularx}
\usepackage{caption}
\usepackage{microtype}
\usepackage[table]{xcolor}
\usepackage{fancyhdr}
\usepackage[colorlinks=true,linkcolor=linkblue,citecolor=linkblue,urlcolor=linkblue]{hyperref}

\definecolor{linkblue}{RGB}{0,80,140}
\definecolor{rulegrey}{RGB}{120,120,120}

\captionsetup{font=small,labelfont=bf,skip=6pt}
\setlength{\parskip}{0pt}
\setlength{\abovedisplayskip}{6pt}
\setlength{\belowdisplayskip}{6pt}

\pagestyle{fancy}
\fancyhf{}
\fancyhead[L]{\footnotesize\itshape Orbital CDN: A Predictive Control Plane for LEO Content Delivery}
\fancyhead[R]{\footnotesize\thepage}
\renewcommand{\headrulewidth}{0.4pt}

% twocolumn defaults to \flushbottom, which forces each column to the page
% bottom by stretching whatever glue it can find. With \parskip at 0pt the only
% stretchable glue is around section headings, so a short column opens a large
% hole beneath its heading. Ragged bottoms are correct here.
\raggedbottom

% --- notation shorthands -----------------------------------------------------
\newcommand{\R}{\mathbb{R}}
\newcommand{\Ind}[1]{\mathbf{1}\!\left[#1\right]}
\newcommand{\ms}{\,\mathrm{ms}}
\newcommand{\km}{\,\mathrm{km}}
\newcommand{\dB}{\,\mathrm{dB}}
\newcommand{\GHz}{\,\mathrm{GHz}}
\newcommand{\mmh}{\,\mathrm{mm/h}}
\newcommand{\Gset}{\mathcal{G}}
\newcommand{\Dset}{\mathcal{D}}
\newcommand{\Cset}{\mathcal{C}}
\newcommand{\Sset}{\mathcal{S}}

\begin{document}

% =============================================================================
\twocolumn[
\begin{@twocolumnfalse}
\begin{center}
{\LARGE\bfseries Orbital CDN: A Physically-Grounded, Predictive\\[2pt]
Control Plane for Content Delivery over LEO Constellations\par}
\vspace{10pt}
{\large Udit Jain\par}
\vspace{4pt}
{\small\ttfamily orbital-cdn.vercel.app \quad$\cdot$\quad github.com/uditjainstjis/orbital-cdn\par}
\vspace{14pt}
\begin{minipage}{0.86\textwidth}
\small
\noindent\textbf{Abstract}---%
Low-Earth-orbit (LEO) constellations are frequently proposed as a substrate for
content delivery on the grounds that the vacuum speed of light exceeds the group
velocity in silica by roughly $47\%$. This argument is necessary but not
sufficient: the ascent and descent legs impose a fixed geometric overhead, so an
orbital path is advantageous only beyond a crossover separation. We present a
routing and control plane that (i) computes end-to-end latency from the actual
geometry of the selected path rather than from a distance surrogate, (ii)
evaluates rain-induced link degradation using a faithful implementation of the
ITU-R P.838-3/P.839-4/P.618-13 recommendation chain, validated against all $28$
published reference coefficients, (iii) expresses every competing routing
objective in a single \emph{latency-equivalent} currency so that policy weights
trade commensurate quantities, and (iv) closes an observe--predict--plan--act--
verify loop in which a gradient-boosted forecaster, trained on four months of
NASA POWER reanalysis precipitation for eight real teleport sites, prices the
probability of imminent link outage into the routing cost.
Under a paired four-arm evaluation over $5{,}600$ requests and $2{,}800$ hours of
recorded weather, request failures fall from $80$ (weather-blind) to $2$
(instantaneous observation) to $0$ (forecast-aware). We report two negative
results that qualify the contribution: the dominant share of the benefit accrues
from \emph{observation} rather than \emph{prediction}, and our hysteretic agent
attains no failure reduction over stateless predictive routing while incurring
$+5.3\ms$ of mean latency. Independently, the simulator reproduces a published
LEO-versus-fibre crossover distance of $4{,}472\km$ without having been
calibrated against it.
\vspace{6pt}

\noindent\textbf{Keywords}---%
LEO constellations, content delivery networks, multi-objective routing, ITU-R
P.618, rain attenuation, probabilistic forecasting, calibration, autonomous
network operations.
\end{minipage}
\end{center}
\vspace{16pt}
\end{@twocolumnfalse}
]

% =============================================================================
\section{Introduction}

The case for routing content through low-Earth orbit rests on a single physical
asymmetry. Light propagates through vacuum at $c_{0}=299{,}792\km\mathrm{/s}$
but through fused-silica fibre at $c_{f}=c_{0}/n\approx204{,}190\km\mathrm{/s}$
for a refractive index $n\approx1.468$, a speed advantage of approximately
$47\%$. A path that ascends to a satellite, traverses a sequence of optical
inter-satellite links (ISLs), and descends to a ground station may therefore
deliver a lower round-trip time (RTT) than a terrestrial path between the same
endpoints.

This argument is incomplete. The ascent and descent legs contribute a
separation-independent overhead of order $2h/\sin\theta$ for shell altitude $h$
and elevation angle $\theta$. Consequently the orbital path is advantageous only
when the terrestrial separation is large enough to amortise that overhead. The
engineering question is therefore not \emph{whether} orbit is faster, but
\emph{where the crossover lies} and \emph{what else is sacrificed} in reaching
it---because a LEO content plane is subject to three constraints with no
terrestrial analogue: orbital eclipse, which removes solar power from a
compute-bearing satellite; transit of the South Atlantic Anomaly (SAA), where the
inner Van Allen belt descends to approximately $200\km$; and rain fade at the
ground segment, which above $10\GHz$ can exceed any economically closable link
margin.

\subsection{Contributions}

\begin{enumerate}
\item \textbf{A standards-faithful attenuation model.} We implement the ITU-R
recommendation chain exactly and validate the implementation against all $28$
published reference coefficients spanning $1$--$30\GHz$
(Section~\ref{sec:physics}). Rain fade given a rain rate is an analytic problem;
we treat it as such rather than approximating it with a learned surrogate.

\item \textbf{A commensurable cost functional.} All routing terms---propagation,
eclipse, radiation, rain fade, learned penalties and predicted outage---are
expressed in latency-equivalent milliseconds, eliminating a class of unit
pathology we encountered and document in Section~\ref{sec:cost}.

\item \textbf{A calibration-first predictive layer.} We show that a persistence
baseline is a strong \emph{discriminator} of imminent fade but a poor
\emph{probability estimator}, and that the learned model's contribution is
almost entirely calibration---the property a cost-weighted router requires
(Section~\ref{sec:predict}).

\item \textbf{A paired evaluation that decomposes the benefit.} By separating an
arm that observes current conditions from one that forecasts, we attribute
$97.5\%$ of the achievable failure reduction to observation and the remainder to
prediction (Section~\ref{sec:results}).

\item \textbf{Reported negative results.} Our hysteretic agent does not improve
on stateless predictive routing in this regime, and we state the conditions
under which it would (Section~\ref{sec:discussion}).
\end{enumerate}

% =============================================================================
\section{Related Work}

Handley~\cite{handley2018} established that a LEO relay network can undercut
terrestrial fibre on long-haul paths, reporting an advantage beyond roughly
$3{,}000\km$ for a $1{,}150\km$ shell. Chaudhry and
Yanikomeroglu~\cite{chaudhry2022} formalise the crossover distance as a function
of shell altitude and fibre refractive index, obtaining $4{,}472\km$ at
$h=550\km$, $n=1.4675$; critically, their terrestrial model assumes great-circle
routing without path inflation, so realised crossovers should be shorter. We use
their figure as an external check rather than a fitted parameter.

Rain attenuation on Earth--space paths is standardised in ITU-R
P.618~\cite{itu618}, with specific attenuation coefficients in
P.838~\cite{itu838} and rain height in P.839~\cite{itu839}. These are the models
used in operational link design; we adopt them unmodified.

Orbital propagation follows the SGP4 analytic theory~\cite{vallado2006} applied
to two-line element sets. Pfandzelter and Bermbach~\cite{pfandzelter2023}
characterise the reliability envelope of in-orbit edge computation, including
single-event-upset rates under SAA transit, and provide the shell parameters we
scale from.

% =============================================================================
\section{System Model and Notation}
\label{sec:model}

Let $\Sset$ denote the satellite set, $|\Sset|=180$, arranged as a Walker-Delta
constellation of $9$ planes $\times\,20$ satellites at inclination $53^{\circ}$
and altitude $h_{s}=550\km$. This is a $1/8$-scale reduction of the authorised
Starlink Gen1 shell~1 ($1{,}584$ satellites, $72\times22$, identical altitude and
inclination), adopted for real-time rendering tractability. A subset
$\Dset\subset\Sset$, $|\Dset|=4$, carries compute at $h_{d}=640\km$. The ground
segment comprises gateways $\Gset$, $|\Gset|=8$, located at published teleport
coordinates, and origin cities $\Cset$, $|\Cset|=8$.

Satellite states are obtained by SGP4 propagation of two-line elements retrieved
from CelesTrak, converted from the ECI frame to geodetic coordinates through
Greenwich Mean Sidereal Time. When the catalogue is unavailable the system falls
back to a closed-form Walker-Delta model,
\begin{align}
\phi(t) &= \arcsin\!\big(\sin i\,\sin\nu(t)\big), \\
\lambda(t) &= \Omega(t) + \operatorname{atan2}\!\big(\cos i\,\sin\nu(t),\,\cos\nu(t)\big),
\end{align}
with inclination $i$, true anomaly $\nu$ advancing at the orbital rate, and right
ascension of the ascending node $\Omega$ advancing at the Earth-rotation rate.

Table~\ref{tab:notation} collects the principal symbols.

\begin{table}[t]
\centering
\caption{Principal notation.}
\label{tab:notation}
\small
\begin{tabularx}{\columnwidth}{@{}lX@{}}
\toprule
Symbol & Meaning \\
\midrule
$g(\cdot,\cdot)$ & great-circle distance (km) \\
$\tau(\cdot,\cdot)$ & round-trip propagation time (ms) \\
$\theta$ & path elevation angle (deg) \\
$\gamma_{R}$ & specific rain attenuation (dB/km) \\
$A_{p}$ & attenuation exceeded for $p\%$ of a year (dB) \\
$M$ & rain-fade link margin (dB) \\
$\Pi$ & instantaneous outage probability \\
$w=(w_{\ell},w_{s},w_{r},w_{x})$ & policy weights \\
$\rho$ & learned penalty in $[0,1]$ \\
$\kappa$ & forecast confidence in $[0,1]$ \\
$\Phi$ & failure cost, latency-equivalent (ms) \\
\bottomrule
\end{tabularx}
\end{table}

% =============================================================================
\section{Link Physics}
\label{sec:physics}

\subsection{Specific attenuation (ITU-R P.838-3)}

The specific attenuation obeys the power law
\begin{equation}
\gamma_{R} = k\,R^{\alpha} \qquad [\mathrm{dB/km}],
\label{eq:powerlaw}
\end{equation}
for rain rate $R$ in $\mmh$. The polarisation-resolved coefficients are fitted as
sums of Gaussians in $\log_{10}f$ with $f$ in GHz. For $k$,
\begin{equation}
\begin{split}
\log_{10}k_{\pi} = \sum_{j=1}^{4} a_{j}
\exp\!\left[-\left(\frac{\log_{10}f-b_{j}}{c_{j}}\right)^{\!2}\right] \\
{}+ m_{k}\log_{10}f + c_{k},
\end{split}
\label{eq:kfit}
\end{equation}
and for $\alpha$, with five terms and \emph{without} the outer exponentiation,
\begin{equation}
\begin{split}
\alpha_{\pi} = \sum_{j=1}^{5} a_{j}
\exp\!\left[-\left(\frac{\log_{10}f-b_{j}}{c_{j}}\right)^{\!2}\right] \\
{}+ m_{\alpha}\log_{10}f + c_{\alpha},
\end{split}
\label{eq:afit}
\end{equation}
where $\pi\in\{H,V\}$ indexes polarisation. The asymmetry between
\eqref{eq:kfit} and \eqref{eq:afit}---$k$ is fitted in logarithmic space and
must be exponentiated, $\alpha$ is not---is a common source of silent
implementation error, since a transposed convention still yields
plausible-magnitude output.

For arbitrary path geometry and polarisation tilt $\tau$ ($45^{\circ}$ for
circular polarisation),
\begin{align}
k &= \tfrac{1}{2}\big[k_{H}+k_{V}+(k_{H}-k_{V})\cos^{2}\theta\cos 2\tau\big],
\label{eq:kcomb}\\
\alpha &= \frac{k_{H}\alpha_{H}+k_{V}\alpha_{V}
+(k_{H}\alpha_{H}-k_{V}\alpha_{V})\cos^{2}\theta\cos 2\tau}{2k}.
\label{eq:acomb}
\end{align}

\paragraph{Validation.} Equations \eqref{eq:kfit}--\eqref{eq:acomb} are evaluated
against the reference table published in the recommendation itself. All $28$
coefficients across $\{1,12,19,20,25,28,30\}\GHz$ reproduce to within $2\%$
(Table~\ref{tab:itu}); the check executes as a unit test, so a transcription
error cannot pass unnoticed.

\begin{table}[t]
\centering
\caption{Validation of the P.838-3 implementation against the published
reference table. All 28 checks pass within 2\%.}
\label{tab:itu}
\small
\begin{tabular}{@{}lrrrr@{}}
\toprule
$f$ (GHz) & $k_{H}$ & $\alpha_{H}$ & $k_{V}$ & $\alpha_{V}$ \\
\midrule
1  & 0.0000259 & 0.9691 & 0.0000308 & 0.8592 \\
12 & 0.02386   & 1.1825 & 0.02455   & 1.1216 \\
19 & 0.08084   & 1.0691 & 0.08642   & 0.9930 \\
20 & 0.09164   & 1.0568 & 0.09611   & 0.9847 \\
25 & 0.1571    & 0.9991 & 0.1533    & 0.9491 \\
28 & 0.2051    & 0.9679 & 0.1964    & 0.9277 \\
30 & 0.2403    & 0.9485 & 0.2291    & 0.9129 \\
\bottomrule
\end{tabular}
\end{table}

The four-order-of-magnitude spread in $k$ between $1\GHz$ and $20\GHz$ is not
merely a validation artefact: it is the quantitative reason that hydrometeor
attenuation is negligible at UHF and dominant at Ka-band, and it disqualifies
UHF-band packet corpora as training data for Ka-band fade prediction. We return
to this in Section~\ref{sec:threats}.

\subsection{Rain height and slant geometry}

Rain height follows P.839-4 as $h_{R}=h_{0}+0.36\km$, where $h_{0}$ is the mean
annual $0^{\circ}$C isotherm height. We use the latitude parameterisation
$h_{0}=5.0-0.075(|\phi|-23)$ km for $|\phi|\ge 23^{\circ}$ and $5.0$ km
otherwise, in place of the gridded map, and flag this as an approximation.

The P.618-13 slant-path procedure then proceeds as follows. For
$\theta\ge5^{\circ}$ the path length below the rain height is
\begin{equation}
L_{s}=\frac{h_{R}-h_{g}}{\sin\theta},\qquad L_{G}=L_{s}\cos\theta,
\end{equation}
with $h_{g}$ the station elevation. The horizontal reduction factor accounts for
the finite spatial extent of a rain cell,
\begin{equation}
r_{0.01}=\left[1+0.78\sqrt{\frac{L_{G}\gamma_{R}}{f}}
-0.38\left(1-e^{-2L_{G}}\right)\right]^{-1},
\label{eq:hred}
\end{equation}
and the vertical adjustment factor, with
$\zeta=\arctan\!\big[(h_{R}-h_{g})/(L_{G}r_{0.01})\big]$ and
$\chi=\max(0,\,36-|\phi|)$,
\begin{equation}
\begin{split}
\nu_{0.01}=\Bigg[1+\sqrt{\sin\theta}\;\times \\
\Bigg(31\big(1-e^{-\theta/(1+\chi)}\big)
\frac{\sqrt{L_{R}\gamma_{R}}}{f^{2}}-0.45\Bigg)\Bigg]^{-1}.
\end{split}
\end{equation}
The effective path length is $L_{E}=L_{R}\nu_{0.01}$ and the attenuation
exceeded for $0.01\%$ of an average year is
\begin{equation}
A_{0.01}=\gamma_{R}L_{E}.
\end{equation}
Extrapolation to other exceedance percentages $p\in[0.001,5]\%$ uses
\begin{equation}
A_{p}=A_{0.01}\left(\frac{p}{0.01}\right)^{-\beta},
\end{equation}
\begin{equation}
\beta = 0.655+0.033\ln p-0.045\ln A_{0.01}-\varsigma(1-p)\sin\theta,
\end{equation}
with $\varsigma$ vanishing for $p\ge1\%$ or $|\phi|\ge36^{\circ}$.

\subsection{Instantaneous fade and outage}

For an observed rain rate the same specific attenuation and horizontal reduction
are applied to the current rate rather than the annual statistic, yielding an
instantaneous fade $A(t)$. We note explicitly that \eqref{eq:powerlaw} is valid
pointwise for any $R$, whereas the P.618 \emph{statistics} presuppose a one
minute integration time; our observations are hourly reanalysis means, which
attenuate peak rates. The resulting fade estimate is therefore
\emph{conservative}, biased against the orbital system rather than for it.

Rather than thresholding fade against margin, we map the residual margin through
a logistic link,
\begin{equation}
\Pi(t)=\sigma\!\left(\frac{A(t)-M}{s}\right)
=\left[1+\exp\!\left(\frac{M-A(t)}{s}\right)\right]^{-1},
\label{eq:outage}
\end{equation}
with $M=6\dB$ and softness $s=1\dB$. The smooth transition encodes estimator
uncertainty: a hard threshold would assert a precision the hourly-mean input
does not support.

% =============================================================================
\section{Path Geometry and Latency}

Great-circle separation uses the haversine form
\begin{equation}
g(a,b)=2R_{\oplus}\arcsin\!\sqrt{\sin^{2}\tfrac{\Delta\phi}{2}
+\cos\phi_{a}\cos\phi_{b}\sin^{2}\tfrac{\Delta\lambda}{2}}.
\end{equation}
Ground-to-satellite legs use the slant range
$\sqrt{g^{2}+h^{2}}$; inter-satellite legs are measured as arcs at orbital
radius, $g\,(R_{\oplus}+h)/R_{\oplus}$.

The complete request path is
\begin{equation}
\underbrace{c \to u}_{\text{slant}} \;\to\;
\underbrace{u \to \cdots \to d}_{\text{ISL arcs}} \;\to\;
\underbrace{d \to u'}_{\text{ISL arc}} \;\to\;
\underbrace{u' \to \eta}_{\text{slant}},
\end{equation}
for origin $c$, uplink satellite $u$, data centre $d\in\Dset$, a relay $u'$ in
view of gateway $\eta\in\Gset$, terminating at the gateway. Two properties of
this construction deserve emphasis. First, the descent is routed via a satellite
\emph{actually in view of the gateway}; a direct $d\to\eta$ slant range would
traverse the planetary interior whenever $d$ lies on the far side. Second, the
chain already egresses and returns to the ground, so its length is the complete
request--response path and must \emph{not} be doubled to obtain RTT. Both are
errors we committed and corrected; the second inflated one origin's path by
$45{,}000\km$.

End-to-end latency is then
\begin{equation}
\mathrm{RTT}=\frac{\ell_{\mathrm{path}}}{c_{0}}\cdot10^{3}
+ T_{\mathrm{proc}} + T_{\mathrm{wx}},
\end{equation}
with compute term $T_{\mathrm{proc}}$ and rain-fade delay $T_{\mathrm{wx}}$.

\subsection{Terrestrial baseline}

The comparison baseline is \emph{measured}, not modelled: ICMP round-trip times
between third-party vantage points, cross-checked against a public cloud
provider's published $30$-day median backbone latencies
(Table~\ref{tab:baseline}). These are datacentre-to-datacentre and therefore
exclude consumer access latency, making the baseline conservative. A physical
model,
\begin{equation}
\mathrm{RTT}_{f}=2\left(\frac{\varrho\, g}{c_{f}}\cdot10^{3}
+ \varsigma_{sw}\frac{\varrho\, g}{10^{3}} + T_{acc}\right)+T_{\mathrm{proc}},
\end{equation}
with route-inflation factor $\varrho=1.42$, switching cost
$\varsigma_{sw}=1\ms$ per $1{,}000\km$ and access latency $T_{acc}=5\ms$, is
retained as a fallback for origins outside the measurement set and is displayed
alongside the measurement so that model error is visible rather than concealed.

\begin{table}[t]
\centering
\caption{Measured terrestrial baselines against the physical model. Distances
are to the nearest terrestrial cloud region.}
\label{tab:baseline}
\small
\begin{tabular}{@{}lrrrr@{}}
\toprule
Origin & $g$ (km) & Measured & Backbone & Model \\
\midrule
New York   &   325 & \textbf{7.5} &   8 &  15 \\
London     &   637 & \textbf{14.4}&  17 &  21 \\
Delhi      & 4{,}145 & \textbf{92.1}&  53 &  79 \\
Dubai      & 4{,}836 & \textbf{122.2}& 100 &  91 \\
Lagos      & 4{,}873 & \textbf{118.5}& --- &  92 \\
Tokyo      & 5{,}322 & \textbf{66.4}&  72 &  99 \\
Sydney     & 6{,}305 & \textbf{92.7}&  95 & 116 \\
S\~{a}o Paulo & 7{,}626 & \textbf{113.1}& 118 & 138 \\
\bottomrule
\end{tabular}
\end{table}

% =============================================================================
\section{Multi-Objective Routing in a Commensurable Currency}
\label{sec:cost}

Data-centre and gateway selection minimise scalarised costs in which
\emph{every} term carries units of latency-equivalent milliseconds:
\begin{align}
C_{\Dset}(c,d) =\;& w_{\ell}\,\tau(c,d)
+ w_{s}\,\kappa_{e}\,\Ind{\mathrm{ecl}(d)} \nonumber\\
&+ w_{r}\,\kappa_{a}\,\Ind{\mathrm{saa}(d)}
+ w_{\ell}\,\kappa_{p}\,\varsigma\,\rho^{\ell}_{d} \nonumber\\
&+ w_{r}\,\kappa_{p}\,\varsigma\,\rho^{r}_{d},
\label{eq:dccost}\\[2pt]
C_{\Gset}(c,\eta) =\;& w_{\ell}\,\tau(c,\eta)
+ w_{x}\,T_{\mathrm{wx}}(\eta) \nonumber\\
&+ w_{x}\,\kappa_{p}\,\varsigma\,\rho^{x}_{\eta}
+ w_{x}\,J(\eta),
\label{eq:gwcost}
\end{align}
with $\tau(a,b)=2g(a,b)/c_{0}\cdot10^{3}$, eclipse cost $\kappa_{e}=25\ms$,
radiation cost $\kappa_{a}=40\ms$, adaptation scale $\kappa_{p}=30\ms$,
adaptation gain $\varsigma=0.6$, and predictive term $J$ defined in
\eqref{eq:predcost}. Selection is $\arg\min$ over each candidate set.

\paragraph{Why commensurability is load-bearing.} An earlier formulation mixed a
distance normalised to $[0,1]$ with indicator penalties of unit magnitude. Under
a balanced weighting the eclipse indicator then dominated half a planetary
circumference of detour, and the router dispatched requests to the antipodal
hemisphere in pursuit of insolation. Stating the eclipse penalty as
``$25\ms$ of detour at unit weight'' renders the trade explicit and auditable;
it also makes the weights interpretable to an operator, which a normalised
formulation does not.

\paragraph{Emergent policy structure.} With four weight vectors the resulting
operating points form a Pareto frontier that was not designed
(Table~\ref{tab:policy}). Notably, the reliability-weighted policy attains
\emph{exactly} zero rain exposure---an unspecified consequence of
$w_{x}=0.90$---at a cost of $12\ms$ in median latency.

\begin{table}[t]
\centering
\caption{Emergent operating points. Rain exposure is the fraction of served
requests whose gateway was in precipitation.}
\label{tab:policy}
\small
\begin{tabular}{@{}lcccc@{}}
\toprule
Policy & $w_{\ell}$ & $w_{s}/w_{r}/w_{x}$ & $p_{50}$ RTT & Rain exp. \\
\midrule
latency  & 0.95 & 0.05/0.05/0.05 & \textbf{61 ms} & 13.7\% \\
balanced & 0.50 & 0.50/0.50/0.40 & 67 ms & 13.9\% \\
reliable & 0.20 & 0.30/0.95/0.90 & 73 ms & \textbf{0.0\%} \\
green    & 0.20 & 0.90/0.20/0.30 & 78 ms & 5.5\% \\
\bottomrule
\end{tabular}
\end{table}

% =============================================================================
\section{Adaptation from Observed Telemetry}

Every served request appends a record to an append-only log. Over a
user-selected window $W$, aggregates induce penalties $\rho\in[0,1]$ that
re-enter \eqref{eq:dccost}--\eqref{eq:gwcost}. For gateways,
\begin{equation}
\rho^{x}_{\eta}=\min\!\left(1,\;
\frac{1}{|W_{\eta}|}\sum_{t\in W_{\eta}} T_{\mathrm{wx}}(\eta,t)\;/\;T^{\max}_{\mathrm{wx}}\right).
\end{equation}
For data centres we separate two signals that carry different policy
coefficients: an observed tail-latency excess relative to the best performer,
\begin{equation}
\rho^{\ell}_{d}=\min\!\left(1,\;
\frac{q_{95}(d)-\min_{d'}q_{95}(d')}{\min_{d'}q_{95}(d')}\right),
\end{equation}
and an observed SAA-transit fraction $\rho^{r}_{d}$.

\paragraph{An information-theoretic exclusion.} We deliberately do \emph{not}
learn a historical eclipse fraction. The cost function already reads
$\Ind{\mathrm{ecl}(d)}$ instantaneously; a windowed average of the same
deterministic quantity is strictly less informative than the indicator it would
accompany, and charging both double-counts identical orbital geometry. The
general principle we adopt is that a learned term is admissible only if it
encodes information not available to the controller by direct observation. Under
this criterion gateway weather qualifies---it is genuinely stochastic---whereas
eclipse, being deterministic celestial mechanics, does not.

\paragraph{Effect.} Toggling adaptation alters the selected route in $6$ of $32$
origin--policy combinations; varying the analysis window alters it in $4/32$
($7$d versus $30$d) and $5/32$ ($24$h versus $7$d).

% =============================================================================
\section{Predictive Layer}
\label{sec:predict}

\subsection{Problem statement}

Let $A_{\eta}(t)$ denote fade at gateway $\eta$. We estimate
\begin{equation}
p_{h}(\eta,t)=\Pr\!\left[\max_{t<t'\le t+h} A_{\eta}(t') > A^{\ast}
\;\middle|\; \mathcal{F}_{t}\right],
\end{equation}
for horizons $h\in\{1,3,6,12\}$ hours and threshold $A^{\ast}=3\dB$, where
$\mathcal{F}_{t}$ is the filtration generated by observations up to $t$.

\subsection{Data and leakage control}

Precipitation is NASA POWER hourly reanalysis for the eight teleport
coordinates, $2{,}928$ hours per site. We note a unit convention that is easy to
mishandle: the hourly endpoint reports \texttt{PRECTOTCORR} in
$\mathrm{mm/day}$ despite hourly sampling, requiring division by $24$; omitting
this inflates every rate by $24\times$ and renders all downstream fade estimates
meaningless.

Twenty-one features are constructed, all measurable at or before $t$: lagged
rates at $\{1,2,3,6,12\}$ hours, first differences over $1$ and $3$ hours,
rolling means over $\{3,6,24\}$ hours, rolling maxima over $\{6,24\}$ hours, a
$24$-hour wet fraction, instantaneous fade, hours since last wet, site
climatology, absolute latitude, and diurnal and seasonal phase.

Partitioning is \emph{temporal}, never random: consecutive hours are strongly
dependent, and a random split leaks future information into training. Training
uses hours $[0,1756)$, validation $[1756,2196)$, and test $[2196,2928)$. A
second partition holds out entire \emph{sites} to assess transfer to a gateway
never observed.

\subsection{Results and the primacy of calibration}

Table~\ref{tab:model} reports held-out performance against three baselines. The
learned model is a gradient-boosted ensemble ($90$ estimators, $12$ leaves,
depth $5$), exported to $71$ KB of JSON and evaluated client-side.

\begin{table}[t]
\centering
\caption{Held-out forecast performance. Persistence dominates on ranking at long
horizons; the learned model dominates on calibration at every horizon.}
\label{tab:model}
\small
\begin{tabular}{@{}lcccc@{}}
\toprule
& \multicolumn{2}{c}{Persistence} & \multicolumn{2}{c}{Gradient boosting} \\
\cmidrule(lr){2-3}\cmidrule(lr){4-5}
Horizon & AUC & Brier & AUC & Brier \\
\midrule
$+1$h  & 0.998 & 0.0294 & \textbf{0.999} & \textbf{0.0030} \\
$+3$h  & 0.990 & 0.0290 & \textbf{0.991} & \textbf{0.0071} \\
$+6$h  & \textbf{0.973} & 0.0307 & 0.949 & \textbf{0.0139} \\
$+12$h & \textbf{0.944} & 0.0329 & 0.892 & \textbf{0.0271} \\
\bottomrule
\end{tabular}
\end{table}

The pattern is unambiguous and, we argue, is the substantive finding of this
section. Persistence is an excellent \emph{ordering} of imminent fade---it is
difficult to beat on AUC, and at $6$ and $12$ hours we do not beat it---but it is
a poor \emph{probability}. Brier score improves by factors of $2$ to $10$ at
every horizon. Because the controller multiplies $p_{h}$ by a cost in
\eqref{eq:predcost} rather than merely sorting by it, calibration and not
discrimination is the operative property. A system that ranks perfectly but
reports $0.9$ where the truth is $0.3$ will systematically over-purchase
latency to avoid failures that were never likely.

Held-out AUC on entirely unseen sites is $0.999$, $0.994$, $0.987$ and $0.958$
across the four horizons, indicating that the learned structure is not merely
site memorisation.

\subsection{Predictive cost and confidence gating}

The predicted-failure term entering \eqref{eq:gwcost} is
\begin{equation}
J(\eta)=\Phi\left[\Pi_{\eta}
+ \kappa(\eta)\,\bar{p}(\eta)\,\big(1-\Pi_{\eta}\big)\right],
\label{eq:predcost}
\end{equation}
where $\bar{p}=\sum_{h}\omega_{h}p_{h}\big/\sum_{h}\omega_{h}$,
$\Phi=900\ms$, and horizon weights
$\omega=(1.0,0.55,0.25,0.10)$. Only the forecast is attenuated by confidence
$\kappa$; the present observation $\Pi_{\eta}$ is certain and enters undamped.

Confidence is reduced when the operating point departs from the training
support---an out-of-distribution rate of $90\mmh$ at a site whose recorded
maximum is $11\mmh$ yields $\kappa=0.45$, below the autonomy threshold
$\kappa_{\min}=0.60$. Below that threshold the agent may recommend but not act.

\paragraph{Degradation policy.} If the predictive subsystem raises, times out, or
is disabled, $J\equiv0$ and the deterministic controller of
Section~\ref{sec:cost} is recovered unchanged. Learned components are never a
single point of failure for routing.

% =============================================================================
\section{The Operations Agent}
\label{sec:agent}

\subsection{Formalisation}

The agent is a deterministic policy engine---explicitly not a learned
controller---executing on a fixed cadence decoupled from rendering. At tick $t$
it observes the network state, evaluates candidate routes, and emits an action
from $\{\textsc{Hold},\textsc{Propose},\textsc{Reroute}\}$.

Define the horizon-weighted forward risk
$\bar{r}(\eta)=\sum_{h}\omega_{h}p_{h}(\eta)\big/\sum_{h}\omega_{h}$, and let
$\eta_{t}$ be the incumbent gateway and
$\eta^{\star}=\arg\min_{\eta}C_{\Gset}(c,\eta)$. The agent reroutes iff the
conjunction
\begin{equation}
\Theta_{t} \;\wedge\;
\Lambda_{t} \;\wedge\;
\Xi_{t} \;\wedge\;
\Delta C_{t}\ge\delta \;\wedge\;
\Delta r_{t}\ge\epsilon \;\wedge\;
\kappa\ge\kappa_{\min}
\label{eq:gate}
\end{equation}
holds, where
\begin{align}
\Theta_{t} &: \;\bar{r}(\eta_{t})\ge r_{\mathrm{trig}} \;\vee\; \Pi_{\eta_{t}}\ge\Pi_{\mathrm{trig}}
&&\text{(trigger)}, \nonumber\\
\Lambda_{t} &: \; t-t_{\mathrm{route}} \ge \Lambda_{\min}
&&\text{(dwell)}, \nonumber\\
\Xi_{t} &: \; t-t_{\mathrm{act}} \ge \Xi_{\min}
&&\text{(cooldown)}, \nonumber\\
\Delta C_{t} &= C_{\Gset}(c,\eta_{t})-C_{\Gset}(c,\eta^{\star})-\varpi
&&\text{(net gain)}, \nonumber\\
\Delta r_{t} &= \bar{r}(\eta_{t})-\bar{r}(\eta^{\star})
&&\text{(risk drop)}. \nonumber
\end{align}
Here $\varpi$ is an explicit switching cost. Operationally
$r_{\mathrm{trig}}=0.45$, $\Pi_{\mathrm{trig}}=0.30$,
$\Lambda_{\min}=20\,$s, $\Xi_{\min}=15\,$s, $\delta=25\ms$, $\epsilon=0.15$,
$\varpi=12\ms$.

\paragraph{Stability.} The dwell and cooldown predicates render the switching
sequence non-Zeno by construction: consecutive actions are separated by at least
$\min(\Lambda_{\min},\Xi_{\min})$, bounding the switching rate independently of
forecast volatility. The margin conditions $\delta,\epsilon$ impose hysteresis,
so a marginal cost inversion cannot induce oscillation between two gateways of
near-equal cost---the failure mode that makes a naive greedy controller worse
than no controller at all.

\subsection{Authority and verification}

Three authority levels are exposed: \textsc{Off} (advisory suppressed),
\textsc{Assist} (propose and await human confirmation), and \textsc{Autopilot}
(act). The default is \textsc{Assist}. This is a deliberate design position
rather than an implementation limitation: no operational satellite programme
presently delegates unsupervised manoeuvre or reconfiguration authority to a
predictive model, and a system claiming otherwise would be dismissed by any
reviewer with domain experience.

Every action is subsequently scored against realised conditions. For a decision
taken at $t$ with stated risk $p$, the realised outcome
$y=\Ind{\max_{t<t'\le t+h}A(t')>A^{\ast}}$ is recorded and the Brier score
$\mathcal{B}=\frac{1}{N}\sum(p_i-y_i)^{2}$ and reliability diagram are
maintained. A forecast that is never scored is an assertion, not a prediction.

% =============================================================================
\section{Experimental Methodology}
\label{sec:method}

We evaluate under a \emph{paired} design: all arms are exposed to an identical
request schedule against an identical weather realisation, replayed from a
common epoch. Without pairing, inter-arm variance would be dominated by which
arm happened to execute during a storm.

Arms are distinguished solely by their observational endowment:
\begin{itemize}
\item[\textbf{A}] \textsc{Blind} --- static weather label only; the
pre-predictive controller.
\item[\textbf{R}] \textsc{Reactive} --- observes instantaneous fade; no forecast.
\item[\textbf{B}] \textsc{Predictive} --- observes fade and forecast.
\item[\textbf{C}] \textsc{Autopilot} --- as B, plus a per-origin standing route
maintained by the agent of Section~\ref{sec:agent}.
\end{itemize}

Isolating \textbf{R} from \textbf{B} is the methodological crux: it separates the
value of \emph{observing} from the value of \emph{forecasting}, a decomposition
that a three-arm design conflates and that determines whether a predictive
subsystem is justified at all.

Outcome is adjudicated by physics identically in every arm: a request fails iff
$A_{\eta}(t)>M$. The schedule comprises $5{,}600$ requests over $2{,}800$
recorded hours, with origins and policies drawn from a fixed pseudorandom
sequence.

% =============================================================================
\section{Results}
\label{sec:results}

\begin{table}[t]
\centering
\caption{Paired four-arm evaluation. $5{,}600$ requests, $2{,}800$ hours of
recorded weather, identical schedule and realisation in every arm.}
\label{tab:arms}
\small
\setlength{\tabcolsep}{4pt}
\begin{tabular}{@{}lrrrr@{}}
\toprule
Arm & Failures & Rate & Mean (ms) & $p_{95}$ (ms) \\
\midrule
A \textsc{Blind}      & 80 & 1.43\% & 4.0 & 12.7 \\
R \textsc{Reactive}   &  2 & 0.04\% & \textbf{3.7} & \textbf{10.0} \\
B \textsc{Predictive} & \textbf{0} & \textbf{0.00\%} & 3.9 & 11.6 \\
C \textsc{Autopilot}  & \textbf{0} & \textbf{0.00\%} & 9.2 & 31.9 \\
\bottomrule
\end{tabular}
\end{table}

Table~\ref{tab:arms} decomposes the benefit. Of $80$ baseline failures,
$78$ ($97.5\%$) are eliminated by instantaneous observation alone; forecasting
eliminates the residual $2$ at a cost of $+0.2\ms$ in mean latency. The agent
attains no further failure reduction and incurs $+5.3\ms$ relative to arm B,
executing $9$ reroutes of which $7$ ($78\%$) preceded actual degradation of the
abandoned gateway.

\paragraph{External validation.} Using measured baselines, the system loses at
London ($637\km$, $0\%$) and New York ($325\km$, $0\%$) and wins from Delhi
($4{,}145\km$, $\approx90\%$) outward, placing its empirical crossover in the
same region as the $4{,}472\km$ reported by~\cite{chaudhry2022} for an identical
shell altitude. No parameter was fitted to that figure. Given that
\cite{chaudhry2022} models terrestrial paths without route inflation whereas our
baseline is measured and therefore inflated, a realised crossover at or slightly
below their value is the theoretically expected outcome.

\paragraph{Environmental bound.} A caveat constrains the achievable effect size:
across the four-month record only one of eight gateways experiences fade
exceeding the $6\dB$ margin, in $19\%$ of its hours (maximum $19.9\dB$); three
others do so in under $0.6\%$ of hours and four never. The measurable benefit of
any weather-aware routing is bounded above by this exposure.

% =============================================================================
\section{Discussion}
\label{sec:discussion}

\subsection{Observation dominates prediction}

The decomposition in Table~\ref{tab:arms} does not support a strong claim for
forecasting in this regime. Reporting only arms A and B would have licensed the
statement ``predictive routing eliminated $100\%$ of failures''---literally true
and materially misleading, since a substantially simpler reactive controller
captures $97.5\%$ of that benefit.

We attribute this to the temporal structure of the hazard. Rain cells at these
sites persist over multiple hours, so the instantaneous state is already highly
informative about the near future; this is the same property that makes
persistence such a strong AUC baseline in Table~\ref{tab:model}. Forecasting
should be expected to contribute materially only where the hazard has short
autocorrelation relative to the actuation latency, or where actuation itself is
slow enough that acting on the present is already too late.

\subsection{When hysteresis is worth its cost}

Arm C is dominated by arm B. The explanation is structural rather than
parametric: in this simulator, re-selecting a gateway per request is costless,
so a standing route with dwell and hysteresis constraints can only retain a
suboptimal assignment between switching opportunities. The agent's machinery
prices a switching cost $\varpi$ that the environment does not actually charge.

This identifies the condition under which the agent becomes justified rather
than merely defensible: environments in which route changes carry genuine
cost---session affinity, cache warmth, key re-establishment, signalling
load---so that $\varpi>0$ in reality and not only in the controller's model. We
regard the present result as correctly characterising the agent's value as
\emph{zero in this regime}, and we report it rather than tuning the environment
until the agent appears necessary.

\subsection{On the admissibility of learned components}

Two decisions in this work were resolved by the same criterion: a learned
component is admissible only where it supplies information unavailable by direct
observation or by an accepted analytic model. Rain fade given a rain rate is
analytic and standardised, so it is computed, not learned. Historical eclipse
fraction is deterministic and already observed, so it is excluded from the
adaptation vector. Future precipitation is genuinely stochastic, so it is
learned. We suggest this criterion generalises: systems that learn quantities
they could compute or observe acquire the appearance of intelligence and the
substance of noise.

% =============================================================================
\section{Threats to Validity}
\label{sec:threats}

\paragraph{Synthetic demand.} Request traffic is simulated and labelled as such
in the interface. Weather, physics and the comparison baselines are not.
Generalisation of the routing decisions to a real demand distribution is
untested.

\paragraph{Reanalysis smoothing.} Hourly means attenuate the sub-hourly peaks
that dominate Ka-band outage. Absolute fade magnitudes are therefore
under-estimated; the direction of the bias opposes our claims.

\paragraph{Scale reduction.} The constellation is $1/8$ of the shell it models.
Inter-satellite link range limits---operationally of order $5{,}400\km$---are not
enforced, which at reduced density permits hops a real system could not close.

\paragraph{Geometric idealisation.} The SAA is approximated by a rectangle over a
region that is elliptical, bifurcating, and drifting westward at approximately
$0.3^{\circ}$/yr. Rain height uses a latitude parameterisation rather than the
gridded isotherm map.

\paragraph{Domain transfer.} We evaluated and rejected a proposal to train the
fade forecaster on a public corpus of UHF ($\approx437$ MHz) satellite packet
receptions. By Table~\ref{tab:itu}, $k$ differs by nearly four orders of
magnitude between that band and Ka; the corpus cannot contain the hydrometeor
signal the model would be required to learn. We record this as a methodological
caution: dataset availability is not evidence of dataset relevance, and the
governing physics should be consulted before a corpus is adopted.

% =============================================================================
\section{Reproducibility}

The physics implementation is validated by an executable test reproducing all
$28$ published coefficients (\texttt{npm run test:itu}). Weather ingestion and
model training are scripted end-to-end
(\texttt{npm run data:weather}, \texttt{npm run train:fade}), with the trained
ensemble exported as inspectable JSON. Held-out metrics, the temporal split
boundaries, feature schema and data provenance are emitted as a machine-readable
artefact alongside the model. The evaluation of Section~\ref{sec:results} is
deterministic under a fixed seed and executes in-browser.

% =============================================================================
\section{Conclusion}

We have described a control plane for LEO content delivery in which the physics
is a faithful standards implementation rather than an approximation, the
comparison baseline is externally measured rather than self-generated, and every
routing objective is expressed in a single commensurable unit. The system
reproduces an independently published crossover distance without having been
fitted to it.

Its principal empirical finding is deflationary and, we believe, more useful for
that: in a regime where the hazard is strongly autocorrelated, most of the
attainable reliability benefit is obtained by \emph{observing} current
conditions, and a predictive subsystem contributes at the margin. The learned
component's value lies in calibration rather than discrimination, and the
hysteretic agent's value in this environment is nil. Establishing the conditions
under which each becomes load-bearing---short hazard autocorrelation for the
forecaster, non-zero switching cost for the agent---is the natural continuation
of this work.

The broader claim we wish to make is not that orbital delivery is faster. It is
that the boundary can be located, that locating it requires physics and
measurement rather than assertion, and that a system which reports where it
fails is more useful than one which reports only where it succeeds.

% =============================================================================
\begin{thebibliography}{9}
\small

\bibitem{handley2018}
M.~Handley,
``Delay is not an option: Low latency routing in space,''
in \emph{Proc. ACM HotNets}, 2018, pp.~85--91.

\bibitem{chaudhry2022}
A.~U.~Chaudhry and H.~Yanikomeroglu,
``When to crossover from Earth to space for lower latency data communications?,''
\emph{arXiv:2203.00154}, 2022.

\bibitem{itu618}
ITU-R,
``Propagation data and prediction methods required for the design of
Earth-space telecommunication systems,''
\emph{Recommendation ITU-R P.618-13}, Geneva, 2017.

\bibitem{itu838}
ITU-R,
``Specific attenuation model for rain for use in prediction methods,''
\emph{Recommendation ITU-R P.838-3}, Geneva, 2005.

\bibitem{itu839}
ITU-R,
``Rain height model for prediction methods,''
\emph{Recommendation ITU-R P.839-4}, Geneva, 2013.

\bibitem{vallado2006}
D.~A.~Vallado, P.~Crawford, R.~Hujsak, and T.~S.~Kelso,
``Revisiting spacetrack report \#3,''
in \emph{Proc. AIAA/AAS Astrodynamics Specialist Conf.}, 2006.

\bibitem{pfandzelter2023}
T.~Pfandzelter and D.~Bermbach,
``Edge computing in low-earth orbit --- what could possibly go wrong?,''
in \emph{Proc. ACM LEO-NET}, 2023.

\bibitem{power}
NASA Langley Research Center,
``POWER: Prediction of Worldwide Energy Resources, hourly API,''
\url{https://power.larc.nasa.gov}, accessed 2026.

\bibitem{fcc2021}
Federal Communications Commission,
``Order and Authorization, IBFS File No. SAT-MOD-20200417-00037,''
FCC 21-48, Apr.~2021.

\end{thebibliography}

\end{document}
